\scanpage{37}
\mspage{11}

En résumé et sauf erreur

\uline{Prop} Soit $X$ une sphère conforme \deleted{\unsure{orientée|marquée}}, $H\subset\mathrm{Aut}_{\mathrm{conf}}(X)\overset{\mathrm{déf}}=G$ un \deleted{\unsure{g}}sous-groupe \uline{compact} \deleted{de}, Alors

a) Il existe une métrique riemannienne sphérique normalisée sur $X$ comp. avec sa structure conforme

b) Celle-ci est unique sss \deleted{il n'existe pas de}\inserted{$X^H=\varnothing$}

\deleted{\unsure{pt}}\deleted{\unsure{b|c}}c) Supposons $\deleted{\unsure{X^H=|X^H\ne\varnothing}}$ $X^H\ne\varnothing$, $H\ne\{1\}$. Alors On a \deleted{\unsure{$X$}}\inserted{ou bien} $\operatorname{card}X^H=2$, ou bien $H=\{1,\sigma\}$ où $\sigma$ est une inversion.

1) Si $\operatorname{card}X^H=2$, \deleted{alors} $H\inserted{{}^+=H\cap G^+.\ \text{On a }X^H=X^{H^+},\ \text{et }H^+}$ est contenu dans un unique tore compact $T$, et \deleted{les \unsure{la structure|métriques sphériques}}\inserted{les métriques} sphériques riemanniennes normalisées comp., invariantes par \deleted{\unsure{$T$}} $H$, sont en corr. biunivoque avec les orbites de \inserted{$T$ dans} \deleted{(}$X\setminus X^H$\deleted{$)/T$} [donc tte orbite de \deleted{\unsure{$X$}} $H^+$ dans $X\setminus X^H$ définit une telle métrique]

% slanted
\margin{\uline{à vérifier} (\uline{NB} Il faudrait arriver à traduire en termes d'alg. linéaire, l'action d'un autom. \deleted{\unsure{de}} $\in G$ sur une métrique sphérique\ldots)}

\deleted{\unsure{d)|1')}} 2) Si $H=\{1,\sigma\}$, $\sigma$ une inversion, les métriques invariantes sous $H$ sont en corr. biunivoque avec les éléments de $(X\setminus X^H)/H=(X\setminus x^\sigma)/\sigma$
% archive sheet number: 39
